\section{Introduction}

Les tests ont pour objectif de vérifier que les fonctionnalités développées correspondent aux besoins et
que les règles métier importantes ne se cassent pas lors des modifications. Dans PharmacieConnect, la
validation combine des tests automatisés côté backend et des tests fonctionnels manuels sur les
interfaces web et mobile.

\section{Stratégie de test}

La stratégie retenue est pragmatique. Les règles les plus sensibles sont testées automatiquement :
authentification, audit, restrictions régionales, analytics et imports. Les parcours d'interface sont
validés manuellement, car ils dépendent de l'affichage, de la navigation, des permissions de localisation
et des données chargées.

\begin{table}[H]
    \centering
    \begin{tabular}{L{4cm}L{9cm}}
        \toprule
        \textbf{Type de test} & \textbf{Objectif} \\
        \midrule
        Tests unitaires backend & Vérifier une règle métier isolée dans un service. \\
        Tests d'intégration backend & Vérifier le comportement entre route, service et base de données de test. \\
        Tests d'import & Contrôler la lecture CSV, les validations et les compteurs de succès/échec. \\
        Tests de sécurité & Vérifier les rôles, la vérification email, les jetons et les accès interdits. \\
        Tests manuels web & Valider les pages d'administration, les formulaires et les tableaux. \\
        Tests manuels mobile & Valider la recherche, la carte, le calendrier, les langues et les paramètres. \\
        \bottomrule
    \end{tabular}
    \caption{Stratégie de validation}
\end{table}

\section{Tests automatisés du backend}

Les tests backend sont écrits avec Pytest. La configuration crée une base SQLite en mémoire pour isoler
les tests de la base réelle du développeur. Cela permet de rejouer les tests sans risque sur les données de
production ou de démonstration.

Les fichiers de test présents couvrent notamment :

\begin{itemize}
    \item \texttt{test\_analytics\_dashboard.py} pour les indicateurs du tableau de bord ;
    \item \texttt{test\_assistant\_regions.py} pour les restrictions régionales des assistants ;
    \item \texttt{test\_audit\_logs.py} pour la traçabilité des connexions et des actions ;
    \item \texttt{test\_medicine\_upload.py} pour l'import et la recherche de médicaments.
\end{itemize}

\subsection{Validation de l'audit}

Les tests vérifient qu'une connexion administrateur réussie crée une entrée
\texttt{ADMIN\_LOGIN}, qu'un échec crée une entrée \texttt{ADMIN\_LOGIN\_FAILED}, et qu'une inscription crée
une entrée \texttt{USER\_REGISTERED}. Cette validation est importante car l'audit est une fonctionnalité de
confiance : s'il manque des événements, l'administrateur ne peut plus suivre correctement l'activité.

\subsection{Validation des restrictions régionales}

Les assistants ne doivent pas gérer toutes les pharmacies du pays. Les tests vérifient que les opérations
acceptent uniquement les gouvernorats appartenant à la région de l'assistant. Par exemple, un assistant
du nord ne doit pas pouvoir importer une pharmacie localisée dans le sud.

\subsection{Validation des imports de médicaments}

Les tests d'import de médicaments vérifient les colonnes obligatoires, la normalisation des prix, la gestion
des doublons par \texttt{code\_pct} et les messages d'erreur pour les lignes invalides. Cette partie est
sensible parce qu'un fichier mal formé peut introduire beaucoup d'erreurs dans le catalogue.

\section{Scénarios de validation fonctionnelle}

Les scénarios suivants décrivent les parcours validés sur les interfaces.

\begin{longtable}{L{4cm}L{5cm}L{5cm}}
    \caption{Scénarios fonctionnels testés}\\
    \toprule
    \textbf{Scénario} & \textbf{Étapes} & \textbf{Résultat attendu} \\
    \midrule
    \endfirsthead
    \toprule
    \textbf{Scénario} & \textbf{Étapes} & \textbf{Résultat attendu} \\
    \midrule
    \endhead
    Connexion admin & Saisir email et mot de passe, valider le formulaire. & Accès au tableau de bord et stockage des jetons. \\
    Recherche mobile & Entrer un nom ou une ville dans l'écran d'accueil. & Liste filtrée de pharmacies avec informations lisibles. \\
    Recherche par proximité & Autoriser la localisation et lancer la recherche. & Pharmacies triées par distance dans le rayon configuré. \\
    Import pharmacies & Charger un fichier CSV valide depuis l'administration. & Nombre de lignes réussies, erreurs éventuelles et données visibles. \\
    Import gardes & Charger un planning de garde CSV. & Gardes consultables par l'admin et par le mobile pour la date donnée. \\
    Recherche médicament & Chercher par nom, DCI ou code PCT. & Résultats pertinents et détail du médicament accessible. \\
    Audit logs & Filtrer par action, entité ou date. & Résultats paginés cohérents avec le compteur. \\
    Changement de langue & Passer du français à l'arabe dans les paramètres. & Interface traduite et direction RTL appliquée. \\
    \bottomrule
\end{longtable}

\section{Validation des API}

La validation de l'API se fait à deux niveaux. D'abord, FastAPI fournit une documentation OpenAPI
permettant de vérifier les routes, les paramètres attendus et les modèles de réponse. Ensuite, les clients
web et mobile consomment les mêmes endpoints, ce qui permet de repérer rapidement une incohérence de
contrat.

Quelques règles de validation sont particulièrement importantes :

\begin{itemize}
    \item les coordonnées GPS doivent rester dans les bornes latitude/longitude ;
    \item les limites de pagination ne doivent pas accepter des valeurs excessives ;
    \item les fichiers CSV sont limités à 5 Mo et 5000 lignes ;
    \item les routes d'administration exigent un compte autorisé ;
    \item les dates de garde doivent être valides avant l'enregistrement.
\end{itemize}

\section{Résultats obtenus}

Les tests automatisés existants couvrent les règles les plus critiques du backend. Les tests manuels ont
confirmé que les parcours principaux sont utilisables : connexion, tableau de bord, gestion des
pharmacies, import de gardes, recherche mobile, carte, calendrier et catalogue des médicaments.

La validation a aussi permis d'identifier des limites. Certaines vérifications restent manuelles, notamment
la qualité visuelle de l'interface mobile, le comportement exact de la géolocalisation selon l'appareil et
les performances sur un volume de données plus grand. Ces points sont acceptables pour une première
version, mais ils doivent être renforcés avant un déploiement large.

\section{Limites de la validation}

La validation réalisée ne remplace pas une campagne de tests en conditions réelles. Les principales
limites sont :

\begin{itemize}
    \item absence de tests automatisés complets côté interface web ;
    \item absence de tests end-to-end couvrant mobile, API et base simultanément ;
    \item données de test encore limitées par rapport à un usage national ;
    \item comportement GPS dépendant de l'appareil, du système et des permissions ;
    \item pas encore de tests de charge poussés sur les endpoints publics.
\end{itemize}

\section{Conclusion}

La phase de tests a permis de sécuriser les fonctionnalités importantes du backend et de valider les
parcours principaux des interfaces. Le projet dispose d'une base de tests utile, surtout pour l'audit, les
imports, les régions et les analytics. Pour une version de production, il serait pertinent d'ajouter des
tests end-to-end, des tests de charge et une validation mobile sur plusieurs appareils.
